\section{Logbook Entry 009: Target-Variable Sweep for Real Allen Neural Memory}
\label{sec:logbook-entry-009-real-allen-target-sweep}

\subsection{Purpose}

This entry records the first target-variable sweep for the real Allen Neuropixels Neural HRSM memory branch. The purpose was to test whether the weak but controlled memory effect observed for \texttt{population\_mean\_rate\_hz} was specific to mean firing rate, or whether retained neural history was stronger for other population-state variables.

The tested targets were:

\begin{verbatim}
population_mean_rate_hz
population_std_rate_hz
active_unit_fraction
population_l2_rate_norm
population_rate_entropy
population_state_speed
\end{verbatim}

The sweep used the same scaled real Allen v1 population-state matrix from session \texttt{715093703}. The matrix contained 3,012 population-state rows across VISp, VISl, LGd, and CA1, with drifting gratings, natural scenes, and spontaneous activity represented.

\subsection{Method}

The script used for the sweep was:

\begin{verbatim}
scripts/allen/13_real_memory_target_sweep_h5py.py
\end{verbatim}

For each target variable, the script repeated the same memory-kernel audit and shuffled-lag control framework used in the previous entry. The model used lags 1 and 2, ridge regularization with \(\alpha = 1.0\), leave-trial-out cross-validation, and 100 shuffled-lag controls per valid region-stimulus group.

The output files were:

\begin{verbatim}
results/real_allen/session_715093703_h5py_v1/
    real_memory_target_sweep_group_summary.csv
    real_memory_target_sweep_target_summary.csv
    real_memory_target_sweep_shuffles.csv
\end{verbatim}

The group summary contained 72 region-stimulus-target rows. The shuffle table contained 4,800 rows.

\subsection{Target-Level Result}

The target-level summary ranked the targets by controlled memory strength, defined primarily by the median observed-minus-shuffled memory gain.

\begin{center}
\begin{tabular}{lrrrr}
\hline
Target & Median observed gain & Median shuffled mean & Median observed minus shuffled & Rank \\
\hline
population\_rate\_entropy & 0.0116346289 & -0.0445391910 & 0.0876368889 & 1 \\
active\_unit\_fraction & 0.0137359607 & -0.0474621228 & 0.0768096637 & 2 \\
population\_mean\_rate\_hz & 0.0004399958 & -0.0366208999 & 0.0524877861 & 3 \\
population\_l2\_rate\_norm & 0.0063723398 & -0.0356808877 & 0.0411592598 & 4 \\
population\_std\_rate\_hz & 0.0024054294 & -0.0362715271 & 0.0302688966 & 5 \\
population\_state\_speed & -0.0502879022 & -0.0436472333 & 0.0052759079 & 6 \\
\hline
\end{tabular}
\end{center}

The strongest controlled memory signal appeared in \texttt{population\_rate\_entropy}, followed by \texttt{active\_unit\_fraction}. This means the retained-history signal is more visible in population organization and recruitment than in raw state speed.

\subsection{Key Finding}

The target sweep changed the interpretation of the memory result. The memory term is not simply a firing-rate effect. The strongest controlled targets were entropy and active-unit fraction. These are closer to HRSM-native quantities because they describe population organization, participation, and distributional structure.

The result therefore supports a cautious statement: in this real Allen session, retained neural history is most visible in the organization and recruitment of the population state, not in the velocity of state motion.

\subsection{Region-Stimulus Result}

The strongest region-stimulus-target effects were concentrated in VISp spontaneous activity. VISp spontaneous appeared at the top of the ranked group table across nearly all tested targets:

\begin{verbatim}
population_mean_rate_hz, VISp spontaneous:
observed_minus_shuffle_mean = 0.179853

active_unit_fraction, VISp spontaneous:
observed_minus_shuffle_mean = 0.166815

population_l2_rate_norm, VISp spontaneous:
observed_minus_shuffle_mean = 0.163846

population_std_rate_hz, VISp spontaneous:
observed_minus_shuffle_mean = 0.156044

population_rate_entropy, VISp spontaneous:
observed_minus_shuffle_mean = 0.155774
\end{verbatim}

This identifies VISp spontaneous activity as the strongest current target for the next phase. VISl spontaneous and LGd spontaneous also showed controlled memory structure, especially for entropy and active-unit fraction.

\subsection{Negative Result}

The weakest target was \texttt{population\_state\_speed}. Its median observed gain was negative, its fraction of observed-positive groups was zero, and its controlled separation from shuffled lag history was very small:

\begin{verbatim}
median_observed_gain = -0.0502879022
fraction_observed_positive = 0.0
median_observed_minus_shuffle = 0.0052759079
median_empirical_p = 0.3762376238
\end{verbatim}

This indicates that, in this run, memory is not primarily expressed as improved prediction of population-state speed.

\subsection{Visual Summary}

The figure script was:

\begin{verbatim}
scripts/allen/14_make_real_memory_target_sweep_figures_h5py.py
\end{verbatim}

It produced four figures:

\begin{verbatim}
real_allen_target_sweep_ranked_controlled_memory.png
real_allen_target_sweep_observed_vs_shuffle.png
real_allen_target_sweep_top_group_effects.png
real_allen_target_region_controlled_memory_heatmap.png
\end{verbatim}

The figure manifest states the visual conclusion directly: controlled memory is strongest for entropy and active-unit fraction, while population-state speed is not supported in this run.

\subsection{Interpretation}

The target sweep supports a more refined Neural HRSM memory interpretation. The memory axis should not be treated as a generic improvement to all target variables. It appears selectively. In the current session, the retained-history contribution is clearest when predicting how distributed or recruited the population state will become, rather than how fast the state moves.

This is consistent with the broader HRSM framing. A neural memory term should not necessarily mean higher firing. It may mean retained organization, previous recruitment structure, or population-level readiness that shapes the next state.

\subsection{Limitations}

This result is based on one Allen session, one selected extraction scale, one lag set, one ridge model, and one shuffled-lag control. The empirical \(p\)-values are uncorrected and based on 100 shuffles per group. The result should therefore be treated as a prioritization result, not a final inferential result.

\subsection{Conclusion}

The target-variable sweep identifies population entropy, active-unit fraction, and VISp spontaneous activity as the most promising targets for the next Neural HRSM memory analysis. The memory term is weak in raw global \(R^2\), but it becomes more structured when tested against shuffled lag history and when the target is population organization rather than state speed.

\subsection{Next Step}

The next phase should focus on spontaneous activity, especially VISp, VISl, and LGd. The recommended next analysis is a spontaneous-block focused extraction with longer continuous windows, followed by the same HRSM projection, target sweep, and shuffled-lag controls.
